% DERIVED from formal_layer.json — read-only, do not edit.
\begin{assumptionv}{ass:source-iid}[Independent source sampling]
The observed source records \(\mathcal S_n\) are independent and identically distributed with common law \(P_{\mathrm S}\), the law of one observed source record:
\[
\mathcal S_n\sim P_{\mathrm S}^{\otimes n}.
\]
\end{assumptionv}

\begin{assumptionv}{ass:target-iid}[Independent target sampling]
The observed target covariate records \(\mathcal T_n\) are independent and identically distributed with common law \(P_{\mathrm T}\), the law of one observed target covariate:
\[
\mathcal T_n\sim P_{\mathrm T}^{\otimes n}.
\]
\end{assumptionv}

\begin{assumptionv}{ass:sample-independence}[Source–target sample independence]
The observed source records \(\mathcal S_n\) and observed target covariate records \(\mathcal T_n\) are independent:
\[
\mathcal S_n\perp\!\!\!\perp\mathcal T_n.
\]
\end{assumptionv}

\begin{definitionv}{synth_1}[The scalar covariate space]
The covariate space is
\[
\mathcal X := [0,1] \subset \mathbb R.
\]
\end{definitionv}

\begin{assumptionv}{ass:source-density-bounds}[Source density and support bounds]
The source covariate density \(f_{\mathrm S}\), with \(c_f\) and \(C_f\) the fixed lower and upper density envelopes, satisfies the following conditions:
\begin{itemize}
\item \textbf{(Envelope constants.)} \(0<c_f<1<C_f\).
\item \textbf{(Full-data law and support.)} \(P_n^F\), the full-data law, is a probability law. Under this law, the covariate \(X\) and both potential outcomes \(Y(0)\) and \(Y(1)\) belong to \([0,1]\) almost surely.
\item \textbf{(Density identity.)} The covariate marginal of \(P_{\mathrm S}\), the law of one observed source record, has density \(f_{\mathrm S}\) with respect to Lebesgue measure restricted to \(\mathcal X=[0,1]\). Thus, for every Borel set \(B\subseteq\mathbb R\),
\[
P_{\mathrm S}(X\in B)=\int_{B\cap[0,1]} f_{\mathrm S}(x)\,dx.
\]
\item \textbf{(Pointwise bounds.)} For every \(x\in\mathcal X\),
\[
c_f\le f_{\mathrm S}(x)\le C_f.
\]
\end{itemize}
\end{assumptionv}

\begin{assumptionv}{ass:target-density-bounds}[Target covariate density bounds]
The target covariate law \(P_{\mathrm T}\), with target density \(f_{\mathrm T}\), satisfies the following conditions on \(\mathcal X=[0,1]\), the covariate domain, with real constants \(c_f\) and \(C_f\), the lower and upper density envelopes:
\begin{itemize}
\item \textbf{(Density identity.)} For every Borel set \(B\subseteq\mathbb R\),
\[
P_{\mathrm T}(B)=\int_{B\cap\mathcal X}\max\{f_{\mathrm T}(x),0\}\,dx.
\]
\item \textbf{(Density envelope.)} For every \(x\in\mathcal X\),
\[
c_f\le f_{\mathrm T}(x)\le C_f.
\]
\item \textbf{(Positive lower envelope.)}
\[
0<c_f.
\]
\end{itemize}
\end{assumptionv}

\begin{assumptionv}{ass:source-density-holder}[Source density Hölder regularity]
The source covariate density \(f_{\mathrm S}\) belongs to \(\mathcal H^{\beta}([0,1],L)\), the class of continuous functions on \([0,1]\) with supremum norm and \(\beta\)-Hölder seminorm bounded separately by \(L\), where \(\beta\) is the common Hölder exponent and \(L\) is the common Hölder radius:
\[
\beta=\frac18,\qquad L>1,\qquad
f_{\mathrm S}\in\mathcal H^{\beta}([0,1],L).
\]
Thus the bounds are
\[
\lVert f_{\mathrm S}\rVert_{\infty}\le L,
\qquad
[f_{\mathrm S}]_{\beta}\le L,
\]
where \([f_{\mathrm S}]_{\beta}\) denotes the \(\beta\)-Hölder seminorm of the source density.
\end{assumptionv}

\begin{assumptionv}{ass:target-density-holder}[Target density Hölder regularity]
The target covariate density \(f_{\mathrm T}\) satisfies the following conditions:
\begin{itemize}
\item \textbf{(Radius.)} The common Hölder radius \(L\) satisfies \(L>1\).
\item \textbf{(Exponent.)} The common Hölder exponent is \(\beta=1/8\).
\item \textbf{(Regularity.)} \(f_{\mathrm T}\in\mathcal H^\beta([0,1],L)\): it is continuous on \([0,1]\) and satisfies
\[
\sup_{x\in[0,1]}|f_{\mathrm T}(x)|\le L,
\qquad
[f_{\mathrm T}]_\beta
=
\sup_{\substack{x,y\in[0,1]\\x\ne y}}
\frac{|f_{\mathrm T}(x)-f_{\mathrm T}(y)|}{|x-y|^\beta}
\le L.
\]
\end{itemize}
\end{assumptionv}

\begin{assumptionv}{ass:propensity-overlap}[Source encouragement propensity overlap]
The source encouragement propensity \(e\) satisfies
\[
\frac14\le e(x)\le\frac34
\qquad\text{for every }x\in\mathcal X.
\]
\end{assumptionv}

\begin{assumptionv}{ass:propensity-holder}[Hölder regularity of the propensity]
The source encouragement propensity \(e\) belongs to \(\mathcal H^\beta([0,1],L)\), with Hölder exponent \(\beta=1/8\) and common Hölder radius \(L>1\); explicitly,
\[
\lVert e\rVert_\infty\le L,
\qquad
[e]_\beta\le L.
\]
\end{assumptionv}

\begin{assumptionv}{ass:arm-means-holder}[Hölder regularity of arm means]
The source conditional arm means \(m_{Az}\), for responses \(A\in\{Y,D\}\) and encouragement arms \(z\in\{0,1\}\), satisfy the following conditions, with common Hölder radius \(L\) and Hölder exponent \(\beta\):
\begin{itemize}
\item \textbf{(Smoothness.)} \(L>1\), \(\beta=1/8\), and each \(m_{Az}\) is continuous on \([0,1]\), with
\[
\lVert m_{Az}\rVert_\infty\le L,
\qquad
[m_{Az}]_\beta\le L.
\]
\item \textbf{(Range.)} For every \(A\), \(z\), and \(x\in[0,1]\),
\[
0\le m_{Az}(x)\le 1.
\]
\item \textbf{(Source arm moments.)} Under the usual measurability conditions, for every \(A\), \(z\), and measurable set \(B\subseteq[0,1]\),
\[
\int_{\{X\in B,\;Z=z\}} A\,dP_{\mathrm S}
=
\int_B f_{\mathrm S}(x)\,
e(x)^z\bigl(1-e(x)\bigr)^{1-z}\,
m_{Az}(x)\,dx,
\]
where \(P_{\mathrm S}\) is the law of an observed source record, \(f_{\mathrm S}\) is the source covariate density, and \(e\) is the source encouragement propensity.
\end{itemize}
\end{assumptionv}

\begin{assumptionv}{ass:receipt-consistency}[Source treatment receipt consistency]
In the source population, observed receipt equals potential receipt under the realized encouragement:
\[
D=D(Z)\qquad\text{almost surely}.
\]
\end{assumptionv}

\begin{assumptionv}{ass:outcome-consistency}[Outcome consistency and boundedness]
In the source population, the observed outcome equals the potential outcome under the realized receipt and lies in the unit interval, both almost surely:
\[
Y=Y(D), \qquad 0\le Y\le 1.
\]
\end{assumptionv}

\begin{assumptionv}{ass:instrument-randomization}[Conditional instrument randomization]
Let \(O\) denote the primitive full-data record and let \(P_n^{F,\mathrm S}=P_n^F(\,\cdot\mid S=1)\) be its source-population law. The joint source assignment law \(\mathbb P_{\mathrm{assign}}\) of \((O,Z)\) is a probability law, and its \(O\)-marginal is \(P_n^{F,\mathrm S}\). The zero extension of the source propensity,
\[
\widetilde e(x)=
\begin{cases}
e(x),&x\in[0,1],\\
0,&x\notin[0,1],
\end{cases}
\]
is measurable. Put \(\pi_1(x)=e(x)\) and \(\pi_0(x)=1-e(x)\). For every measurable full-data event \(B\) and each \(z\in\{0,1\}\),
\[
\mathbb P_{\mathrm{assign}}(O\in B, Z=z)
=\int_B \pi_z(X(O))\,dP_n^{F,\mathrm S}(O).
\]
Thus encouragement \(Z\) is conditionally independent of \((D(0),D(1),Y(0),Y(1))\) given \(X\) and source membership \(S=1\):
\[
Z\perp\!\!\!\perp (D(0),D(1),Y(0),Y(1))\mid X,S=1.
\]
\end{assumptionv}

\begin{assumptionv}{ass:exclusion}[Exclusion through treatment receipt]
For each \(z\in\{0,1\}\), the source potential outcome generated by setting encouragement to \(z\) equals \(Y(D(z))\) almost surely under \(P_n^F\), the primitive full-data law at sample-size index \(n\). Here \(D(z)\) is potential receipt under encouragement \(z\), and \(Y(d)\) is the potential outcome under receipt \(d\).
\end{assumptionv}

\begin{assumptionv}{ass:monotonicity}[Monotonicity of treatment receipt]
The potential receipts \(D(1)\) and \(D(0)\), under unit and zero encouragement respectively, satisfy
\[
D(1)\ge D(0)
\]
almost surely under \(P_n^F\), the primitive full-data law at sample-size index \(n\).
\end{assumptionv}

\begin{assumptionv}{ass:transport-complier-outcome}[Transport of complier outcomes]
Let \(C\) be the indicator of the complier principal stratum, let \(Y(d)\) be the potential outcome under receipt \(d\), and let \(P_n^F\) be the primitive full-data law at sample-size index \(n\). For \(P_{\mathrm T}\)-almost every covariate value \(x\), where \(P_{\mathrm T}\) is the law of one observed target covariate,
\[
\mathbb E_{P_n^F}[(Y(1)-Y(0))C\mid X=x,S=0]
=
\mathbb E_{P_n^F}[(Y(1)-Y(0))C\mid X=x,S=1],
\]
where \(S=0\) denotes target membership and \(S=1\) denotes source membership.
\end{assumptionv}

\begin{assumptionv}{ass:transport-complier-share}[Transport of complier shares]
Let \(C\) be the indicator of the complier principal stratum, and let \(P_n^F\) be the primitive full-data law at sample-size index \(n\). For \(P_{\mathrm T}\)-almost every covariate value \(x\), where \(P_{\mathrm T}\) is the law of one observed target covariate,
\[
\mathbb E_{P_n^F}[C\mid X=x,S=0]
=
\mathbb E_{P_n^F}[C\mid X=x,S=1],
\]
where \(S=0\) denotes target membership and \(S=1\) denotes source membership.
\end{assumptionv}

\begin{assumptionv}{ass:positive-strength}[Positive transported first stage]
The transported compliance share and first stage, \(\mu\), satisfies
\[
\mu>0.
\]
\end{assumptionv}

\begin{definitionv}{def:model-class}[Model class $\mathcal M_n$]
The model class \(\mathcal M_n\) consists of law triples
\[
P=(P_n^F,P_{\mathrm S},P_{\mathrm T}),
\]
where \(P_n^F\) is the primitive full-data law, \(P_{\mathrm S}\) is the law of one observed source record, and \(P_{\mathrm T}\) is the law of one observed target covariate, satisfying all of the following restrictions at sample-size index \(n\):
\begin{itemize}
\item \textbf{(Sampling.)} \cref{obj:ass:source-iid,obj:ass:target-iid,obj:ass:sample-independence}.
\item \textbf{(Covariate densities.)} \cref{obj:ass:source-density-bounds,obj:ass:target-density-bounds,obj:ass:source-density-holder,obj:ass:target-density-holder}.
\item \textbf{(Propensity and arm means.)} \cref{obj:ass:propensity-overlap,obj:ass:propensity-holder,obj:ass:arm-means-holder}.
\item \textbf{(Causal restrictions.)} \cref{obj:ass:receipt-consistency,obj:ass:outcome-consistency,obj:ass:instrument-randomization,obj:ass:exclusion,obj:ass:monotonicity}.
\item \textbf{(Transport and strength.)} \cref{obj:ass:transport-complier-outcome,obj:ass:transport-complier-share,obj:ass:positive-strength}.
\end{itemize}
\end{definitionv}

\begin{definitionv}{def:honest-intervals}[Honest interval class \(\mathfrak H_n\)]
The class of uniformly honest connected intervals is
\[
\mathfrak H_n=
\left\{
C_n:
C_n\subseteq\Theta\ \text{is a measurable connected interval},
\quad
\inf_{P\in\mathcal M_n}P(\theta\in C_n)\ge 1-\alpha
\right\}.
\]
Here \(\Theta\) is the bounded domain of the target-population complier effect \(\theta\), \(\mathcal M_n\) is the model class, and \(\alpha\) is the prescribed noncoverage probability.
\end{definitionv}

\begin{definitionv}{def:strength-slice}[First-stage strength slice]
The first-stage strength slice, with \(\mu(P)\) denoting the transported first stage under \(P\), is defined for real parameters \(c_f,C_f,L,a\) and every \(n\in\mathbb N\) by
\[
\mathcal M_n(a)
=
\left\{
P\in\mathcal M_n:
0<a\le\frac14,\quad
a\le\mu(P)\le\frac14
\right\}.
\]
Here \(\mathcal M_n\) is the model class of \cref{obj:def:model-class} with parameters \(c_f,C_f,L\).
\end{definitionv}

\begin{definitionv}{def:length-frontier}[Minimax length frontier \(L_n(a)\)]
The minimax expected-length frontier is
\[
L_n(a)
=
\inf_{C_n\in\mathfrak H_n}
\sup_{P\in\mathcal M_n(a)}
\mathbb E_P\!\left[\lambda_{\Theta}(C_n)\right],
\qquad n\ge n_0,\quad 0<a\le\frac14,
\]
where \(\mathfrak H_n\) is the class of uniformly honest connected intervals, \(\mathcal M_n(a)\) is the model slice with transported first stage between \(a\) and \(1/4\), and \(\lambda_{\Theta}\) denotes Lebesgue length restricted to the effect domain \(\Theta\).
\end{definitionv}

\begin{definitionv}{def:marked-density-vector}[Marked-density vector \(F\)]
The marked-density vector for the model in \cref{obj:def:model-class} is
\[
F(x)=\bigl(f_{\mathrm T}(x),q_0(x),q_1(x),
r_{Y0}(x),r_{Y1}(x),r_{D0}(x),r_{D1}(x)\bigr),
\]
where the source arm densities and response-marked arm densities are
\[
q_z(x)=f_{\mathrm S}(x)P_{\mathrm S}(Z=z\mid X=x),
\qquad
r_{Az}(x)=q_z(x)m_{Az}(x),
\qquad A\in\{Y,D\},\quad z\in\{0,1\}.
\]
Here \(m_{Az}(x)\) is the source conditional mean of response \(A\) in encouragement arm \(z\).
\end{definitionv}

\begin{definitionv}{def:smooth-functional}[Transported contrast functional \(T_A\)]
The target-averaged source contrast for response \(A\in\{Y,D\}\) is
\[
T_A=\int_0^1\Phi_A(F(x))\,dx,
\qquad
\Phi_A(F)=f_{\mathrm T}
\left(\frac{r_{A1}}{q_1}-\frac{r_{A0}}{q_0}\right).
\]
\end{definitionv}

\begin{definitionv}{synth_10}[Calibrated score acceptance set]
The calibrated affine-score acceptance set \(\mathcal A_n\) is defined, for real inputs \(\alpha,c_f,C_f,L\), a sample size \(n\in\mathbb N\), and a pair of source and target samples each containing \(n\) records, by
\[
\mathcal A_n
=
\left\{
\theta\in[-1,1]:
\left|\widehat T_Y-\theta\widehat T_D\right|
\le
2\sqrt{\frac{2C_{\mathrm{mse}}}{\alpha}}\,
n^{-1/3}
\right\}.
\]
Here \(\widehat T_Y\) and \(\widehat T_D\) are the cubic estimators computed from the supplied samples with inputs \(c_f,C_f\), and \(C_{\mathrm{mse}}\) is the mean-square envelope evaluated at \(c_f,C_f,L\). Each cubic estimator equals zero when \(n<n_0\), the estimator threshold, and otherwise equals its pilot integral plus its linear, quadratic, and cubic corrections. The square root and negative power use their total real-valued conventions, including at zero.
\end{definitionv}

\begin{definitionv}{def:score-interval}[Inverted score interval \(I_n\)]
The inverted score interval is
\[
I_n=
\begin{cases}
\Theta,& n<n_0,\\
\mathcal A_n,& n\ge n_0\ \text{and}\ \mathcal A_n\ne\varnothing,\\
\{0\},& n\ge n_0\ \text{and}\ \mathcal A_n=\varnothing.
\end{cases}
\]
Here \(\mathcal A_n\) is the calibrated affine-score acceptance set, \(\Theta\) is the bounded effect domain, and \(n_0\) is the threshold for using the cubic estimator.
\end{definitionv}

\begin{theoremv}{thm:sharp-length-frontier-full-elbow}[Sharp length frontier with elbow]
Fix real constants \(\alpha,c_f,C_f,L\) satisfying:
\begin{itemize}
\item \textbf{(Noncoverage probability.)} \(0<\alpha<1\).
\item \textbf{(Lower density envelope.)} \(0<c_f<1\).
\item \textbf{(Upper density envelope.)} \(1<C_f\).
\item \textbf{(Hölder radius.)} \(1<L\).
\end{itemize}
Let \(I_n\) be the score interval of \cref{obj:def:score-interval}, and let \(L_n(a)\) be the minimax expected restricted length defined in \cref{obj:def:length-frontier}. Set
\[
n_0=256.
\]
There exist finite constants \(0<c_0<C_0\), depending only on \((\alpha,c_f,C_f,L)\), such that \(I_n\in\mathfrak H_n\), the class of honest intervals in \cref{obj:def:honest-intervals}, for every integer \(n>0\), and, for every integer \(n\ge n_0\) and every real \(a\) with \(0<a\le 1/4\),
\[
c_0\min\{1,n^{-1/3}/a\}
\le L_n(a)
\le C_0\min\{1,n^{-1/3}/a\}.
\]
Moreover, for every integer \(n\ge n_0\) and every law \(P\in\mathcal M_n\) as defined in \cref{obj:def:model-class}, writing
\[
\mu(P)=T_D(P)
\]
for the transported first stage, the same interval sequence satisfies
\[
\mathbb E_P\!\left[\lambda_{\Theta}(I_n)\right]
\le C_0\min\{1,n^{-1/3}/\mu(P)\},
\]
where \(\lambda_{\Theta}\) denotes Lebesgue length restricted to the effect domain \(\Theta\).
\end{theoremv}

\begin{definitionv}{def:unequal-sample-handle}[Unequal-sample frontier \(R_{N_{\mathrm S},N_{\mathrm T}}(a)\)]
The unequal-sample minimax expected-length frontier is
\[
R_{N_{\mathrm S},N_{\mathrm T}}(a)
=
\inf_{C_{N_{\mathrm S},N_{\mathrm T}}\in
\mathfrak H^{\mathrm u}_{N_{\mathrm S},N_{\mathrm T}}}
\sup_{P\in\mathcal M^{\mathrm u}_{N_{\mathrm S},N_{\mathrm T}}(a)}
\mathbb E_P\!\left[
\lambda_{\Theta}\!\left(C_{N_{\mathrm S},N_{\mathrm T}}\right)
\right],
\qquad
0<a\le\frac14,
\]
for positive integers \(N_{\mathrm S}\) and \(N_{\mathrm T}\), where \(\lambda_{\Theta}\) denotes Lebesgue length restricted to the effect domain \(\Theta\).

Here \(\mathcal M^{\mathrm u}_{N_{\mathrm S},N_{\mathrm T}}\) is the class of full-data and one-record observed-law triples satisfying every law-level density, smoothness, overlap, causal, transport, and positive-strength condition of the equal-sample model class \(\mathcal M_n\), observed under
\[
P_{\mathrm S}^{\otimes N_{\mathrm S}}
\otimes
P_{\mathrm T}^{\otimes N_{\mathrm T}},
\]
where \(P_{\mathrm S}\) is the law of one observed source record and \(P_{\mathrm T}\) is the law of one observed target covariate. The independent source and target samples contain \(N_{\mathrm S}\) records \((X,Z,D,Y)\) and \(N_{\mathrm T}\) records \(X\), respectively. With \(\mu(P)\) denoting the transported first stage, the strength slice is
\[
\mathcal M^{\mathrm u}_{N_{\mathrm S},N_{\mathrm T}}(a)
=
\left\{
P\in\mathcal M^{\mathrm u}_{N_{\mathrm S},N_{\mathrm T}}:
a\le\mu(P)\le\frac14
\right\}.
\]

Let \(C_{N_{\mathrm S},N_{\mathrm T}}\) range over measurable connected random intervals contained in \(\Theta\) and based on this product experiment. With \(\theta(P)\) denoting the target-population average treatment effect among compliers and \(\alpha\) the prescribed noncoverage probability, define
\[
\mathfrak H^{\mathrm u}_{N_{\mathrm S},N_{\mathrm T}}
=
\left\{
C_{N_{\mathrm S},N_{\mathrm T}}:
\inf_{P\in\mathcal M^{\mathrm u}_{N_{\mathrm S},N_{\mathrm T}}}
P\!\left(
\theta(P)\in C_{N_{\mathrm S},N_{\mathrm T}}
\right)
\ge1-\alpha
\right\}.
\]
\end{definitionv}

\begin{remarkv}{oeq:unequal-sample-frontier}[Unequal-sample length phases]
What is the exact piecewise order of \(R_{N_{\mathrm S},N_{\mathrm T}}(a)\), the unequal-sample minimax expected-length frontier, when the source sample size \(N_{\mathrm S}\) and target sample size \(N_{\mathrm T}\) diverge at unrelated rates? Can a single finite-sample honest interval attain this order in every phase, together with a matching legal instrumental-variable lower family satisfying the no-defiers restriction?
\end{remarkv}

\begin{lemmav}{lem:transported-cace-identification}[Transported complier effect identification]
Let \(n\in\mathbb N\) and let \(P\) be a transport law satisfying:
\begin{itemize}
\item \textbf{(Lower envelope.)} \(0<c_f<1\), where \(c_f\) is the lower density envelope.
\item \textbf{(Upper envelope.)} \(1<C_f\), where \(C_f\) is the upper density envelope.
\item \textbf{(Smoothness radius.)} \(1<L\), where \(L\) is the common Hölder radius.
\item \textbf{(Model membership.)} \(P\in\mathcal M_n\), the model class of \cref{obj:def:model-class} with constants \(c_f,C_f,L\).
\end{itemize}
Write \(P_n^F\) for the full-data law, \(D(z)\) for potential receipt under encouragement \(z\), and \(Y(d)\) for potential outcome under receipt \(d\). For each \(z\in\{0,1\}\), there are versions of the source conditional means such that, almost everywhere under the source covariate law,
\[
m_{Dz}(x)
 =\mathbb E_{P_n^F}[D(z)\mid X=x,S=1],
\qquad
m_{Yz}(x)
 =\mathbb E_{P_n^F}[Y(D(z))\mid X=x,S=1].
\]
Here \(m_{Dz}\) and \(m_{Yz}\) are the source receipt and outcome arm means.

Define the source arm contrasts and the complier indicator by
\[
\Delta_D(x)=m_{D1}(x)-m_{D0}(x),
\qquad
\Delta_Y(x)=m_{Y1}(x)-m_{Y0}(x),
\qquad
C=\mathbf 1\{D(1)=1,\ D(0)=0\}.
\]
There are common versions of the conditional means in each of the following chains for which, for \(P_{\mathrm T}\)-almost every \(x\), where \(P_{\mathrm T}\) is the target covariate law,
\[
\begin{aligned}
\Delta_D(x)
&=\mathbb E_{P_n^F}[D(1)-D(0)\mid X=x,S=1]\\
&=\mathbb E_{P_n^F}[C\mid X=x,S=1]
 =\mathbb E_{P_n^F}[C\mid X=x,S=0],\\
\Delta_Y(x)
&=\mathbb E_{P_n^F}[Y(D(1))-Y(D(0))\mid X=x,S=1]\\
&=\mathbb E_{P_n^F}[(Y(1)-Y(0))C\mid X=x,S=1]\\
&=\mathbb E_{P_n^F}[(Y(1)-Y(0))C\mid X=x,S=0].
\end{aligned}
\]
Moreover,
\[
D(1)-D(0)=C
\qquad P_n^F\text{-almost surely}.
\]
For \(A\in\{D,Y\}\), define the transported contrast using the target covariate density \(f_{\mathrm T}\) by
\[
T_A=\int \Delta_A(x)f_{\mathrm T}(x)\,dx.
\]
The transported first stage equals the target complier share, and the transported outcome contrast equals the target complier-weighted outcome effect:
\[
T_D=\mu=P_n^F(C=1\mid S=0),
\qquad
T_Y=\mathbb E_{P_n^F}[(Y(1)-Y(0))C\mid S=0].
\]
The target complier average treatment effect \(\theta\) satisfies
\[
\theta
=\frac{\mathbb E_{P_n^F}[(Y(1)-Y(0))C\mid S=0]}
       {P_n^F(C=1\mid S=0)}
=\frac{T_Y}{T_D}
\in\Theta,
\]
where \(\Theta\) is the effect domain specified by the model.
\end{lemmav}

\begin{definitionv}{synth_3}[Exact block size]
The block size \(m_b\), for a nonnegative integer \(n\) and an index \(b\in\{0,1,2,3\}\), is defined by
\[
m_b := \max\left\{
\left\lfloor \frac{(b+1)n}{4} \right\rfloor
-
\left\lfloor \frac{bn}{4} \right\rfloor,
0
\right\}.
\]
\end{definitionv}

\begin{definitionv}{synth_2}[Deterministic sample blocks]
The deterministic block \(\mathcal B_b^G\) in sample channel \(G\), for sample size \(n\in\{0,1,2,\ldots\}\) and block index \(b\in\{0,1,2,3\}\), uses zero-based record indices and is defined by
\[
\mathcal B_b^G
=
\left\{
i\in\mathbb Z:
0\le i<n,\quad
\left\lfloor\frac{bn}{4}\right\rfloor
\le i<
\left\lfloor\frac{(b+1)n}{4}\right\rfloor
\right\}.
\]
\end{definitionv}

\begin{definitionv}{synth_19}[Target-density coordinate]
In the marked-density vector, \(t\) denotes the target covariate density:
\[
t:=f_{\mathrm T},\qquad t(x)=f_{\mathrm T}(x),\quad x\in[0,1].
\]
Thus \(F_0=t=f_{\mathrm T}\), and the named-coordinate subscript \(t\) denotes coordinate \(0\): \(V_t:=V_0\) for \(V=F\), \(V=H_K^{(b)}\), \(V=R_K^{(b)}\), or \(V=W_r\). In particular, \(F_t=f_{\mathrm T}\). In \(\operatorname{bump}(t)\) and \(\int_{\Theta}P(t\in C_n)\,dt\), \(t\) is instead a local real argument or integration variable, respectively.
\end{definitionv}

\begin{definitionv}{synth_14}[Coordinate marks \(W_{i,r}\) and sample blocks \(\mathcal B_b^G\)]
The coordinate index set and coordinate ordering are
\[
\mathcal I=\{0,\ldots,6\},
\qquad
(F_0,F_1,F_2,F_3,F_4,F_5,F_6)
=(t,q_0,q_1,r_{Y0},r_{Y1},r_{D0},r_{D1}).
\]
The sample channel of coordinate \(i\in\mathcal I\) is
\[
g(i)=
\begin{cases}
\mathrm T,&i=0,\\
\mathrm S,&i\in\{1,\ldots,6\}.
\end{cases}
\]
Vector, histogram, residual, and mark subscripts follow the correspondence
\[
\begin{array}{c|ccccccc}
i&0&1&2&3&4&5&6\\ \hline
\text{coordinate}&t&q_0&q_1&r_{Y0}&r_{Y1}&r_{D0}&r_{D1}.
\end{array}
\]
For \(V=F\), \(V=H_K^{(b)}\), \(V=R_K^{(b)}\), or \(V=W_r\), the named-coordinate subscripts satisfy
\[
(V_t,V_{q_0},V_{q_1},V_{r_{Y0}},V_{r_{Y1}},V_{r_{D0}},V_{r_{D1}})
=(V_0,V_1,V_2,V_3,V_4,V_5,V_6).
\]

Index the source records and target covariate records from zero:
\[
O_r^{\mathrm S}=(X_r^{\mathrm S},Z_r,D_r,Y_r),
\qquad
X_r^{\mathrm T},
\qquad r=0,\ldots,n-1.
\]
Let \(Y_r^{\mathrm{clip}}\) denote the clipped observable outcome used in
\cref{obj:synth_7,obj:synth_8,obj:def:cubic-estimator}.
The observable coordinate marks, with this same clipping convention throughout, are
\[
\begin{aligned}
W_{0,r}&=1,\\
W_{1,r}&=\mathbf 1\{Z_r=0\},
&
W_{2,r}&=\mathbf 1\{Z_r=1\},\\
W_{3,r}&=\mathbf 1\{Z_r=0\}Y_r^{\mathrm{clip}},
&
W_{4,r}&=\mathbf 1\{Z_r=1\}Y_r^{\mathrm{clip}},\\
W_{5,r}&=\mathbf 1\{Z_r=0\}D_r,
&
W_{6,r}&=\mathbf 1\{Z_r=1\}D_r.
\end{aligned}
\]
The mark \(W_{0,r}\) accompanies target record \(r\); the remaining marks accompany source record \(r\).

For \(G\in\{\mathrm S,\mathrm T\}\) and \(b\in\{0,1,2,3\}\), the blocks and their sizes are
\[
\mathcal B_b^G
=
\left\{
r\in\{0,\ldots,n-1\}:
\left\lfloor\frac{bn}{4}\right\rfloor
\le r<
\left\lfloor\frac{(b+1)n}{4}\right\rfloor
\right\},
\qquad
m_b=\left\lfloor\frac{(b+1)n}{4}\right\rfloor
-\left\lfloor\frac{bn}{4}\right\rfloor.
\]
Their one-based counterparts are
\[
\widetilde{\mathcal B}_b^G
=\{r+1:r\in\mathcal B_b^G\}
=\{\lfloor bn/4\rfloor+1,\ldots,\lfloor(b+1)n/4\rfloor\}.
\]
One-based covariates and marks are defined together by
\[
\widetilde X_s^G=X_{s-1}^G,
\qquad
\widetilde W_{i,s}=W_{i,s-1},
\qquad s=1,\ldots,n.
\]
Thus, for a covariate set \(E\), the corresponding block sums satisfy
\[
\sum_{s\in\widetilde{\mathcal B}_b^{g(i)}}
\widetilde W_{i,s}\,
\mathbf 1\{\widetilde X_s^{g(i)}\in E\}
=
\sum_{r\in\mathcal B_b^{g(i)}}
W_{i,r}\,\mathbf 1\{X_r^{g(i)}\in E\}.
\]
Block \(0\) supplies the training records, and blocks \(1,2,3\) supply the evaluation records. Each coordinate-\(i\) sum uses covariates and marks with the same record index in channel \(g(i)\).
\end{definitionv}

\begin{definitionv}{synth_7}[Blockwise marked histogram definition]
The marked histogram \(H_{i,K}^{(b)}(x)\), for coordinate \(i\), resolution \(K\), block \(b\), and evaluation point \(x\), is defined by
\[
H_{i,K}^{(b)}(x)=
\begin{cases}
0,&n<256,\\[2mm]
\displaystyle
\frac{K}{m_b}
\sum_{\ell=0}^{K-1}
\mathbf 1\{x\in E_{\ell,K}\}
\sum_{r\in\mathcal B_b^{g(i)}}
W_{i,r}\,
\mathbf 1\{X_r^{g(i)}\in E_{\ell,K}\},
&n\ge256.
\end{cases}
\]
Here \(\mathcal S_n\), the sample of source records, and \(\mathcal T_n\), the sample of target covariates, each contain \(n\) records, and the parameters range over
\[
n,K\in\{0,1,\ldots\},\qquad
i\in\mathcal I=\{0,\ldots,6\},\qquad
b\in\{0,\ldots,3\},\qquad
x\in\mathbb R.
\]
For \(K\ge1\), the cells \(E_{\ell,K}\) are
\[
E_{\ell,K}=
\begin{cases}
[\ell/K,(\ell+1)/K),&0\le\ell<K-1,\\
[\ell/K,1],&\ell=K-1.
\end{cases}
\]
For \(K=0\), the cell sum is empty and \(H_{i,0}^{(b)}(x)=0\).

Both samples use zero-based record indices \(r\in\{0,\ldots,n-1\}\), with an empty index set when \(n=0\). For sample channel \(G\in\{\mathrm S,\mathrm T\}\), the deterministic block \(\mathcal B_b^G\) and its size \(m_b\) are
\[
\mathcal B_b^G
=
\left\{
r\in\{0,\ldots,n-1\}:
\left\lfloor\frac{bn}{4}\right\rfloor
\le r<
\left\lfloor\frac{(b+1)n}{4}\right\rfloor
\right\},
\qquad
m_b=
\left\lfloor\frac{(b+1)n}{4}\right\rfloor
-
\left\lfloor\frac{bn}{4}\right\rfloor.
\]
The sample channel \(g(i)\) assigned to coordinate \(i\) is
\[
g(i)=
\begin{cases}
\mathrm T,&i=0,\\
\mathrm S,&i\in\{1,\ldots,6\}.
\end{cases}
\]
The covariate \(X_r^{g(i)}\) is the observed covariate of zero-based record \(r\) in that channel.

For zero-based source record \(r\), let \(Z_r\in\{0,1\}\) denote encouragement, \(D_r\in\{0,1\}\) treatment receipt, and \(Y_r\in\mathbb R\) the observed outcome. The observable marks \(W_{i,r}\) are
\[
\begin{aligned}
W_{0,r}&=1,\\
W_{1,r}&=\mathbf 1\{Z_r=0\},
&
W_{2,r}&=\mathbf 1\{Z_r=1\},\\
W_{3,r}&=\mathbf 1\{Z_r=0\}\max\{0,\min\{1,Y_r\}\},
&
W_{4,r}&=\mathbf 1\{Z_r=1\}\max\{0,\min\{1,Y_r\}\},\\
W_{5,r}&=\mathbf 1\{Z_r=0\}\mathbf 1\{D_r=1\},
&
W_{6,r}&=\mathbf 1\{Z_r=1\}\mathbf 1\{D_r=1\}.
\end{aligned}
\]
The constant mark \(W_{0,r}\) accompanies zero-based target record \(r\); each remaining mark accompanies zero-based source record \(r\). In every block sum, the mark and covariate use the same record index.
\end{definitionv}

\begin{definitionv}{synth_4}[Dyadic pilot resolution rule]
The pilot resolution \(K_0(n)\), for sample size \(n\in\mathbb N\), is defined by
\[
K_0(n)=
\begin{cases}
1, & n<n_0,\\[2pt]
2^{\left\lfloor \frac45\,\frac{\log m_0}{\log 2}\right\rfloor},
& n\ge n_0,
\end{cases}
\]
where \(n_0\) is the threshold for using the cubic estimator and \(m_0\) is the size of block \(0\) at sample size \(n\).
\end{definitionv}

\begin{definitionv}{synth_8}[Clipped marked-density pilot vector]
The pilot vector \(\widehat F(x)\), indexed by \(i\in\{0,\ldots,6\}\), is defined for real envelope parameters \(c_f,C_f\), a sample-size index \(n\in\mathbb N\), source and target samples \((\mathcal S_n,\mathcal T_n)\) each of size \(n\), and \(x\in\mathbb R\). For \(n<256\), set
\[
\widehat F(x)=\left(1,\frac12,\frac12,0,0,0,0\right).
\]
For \(n\ge256\), use \(K_0(n)\), the dyadic pilot resolution with exponent \(4/5\), and the block-\(0\) marked histogram \(H_{i,K_0(n)}^{(0)}(x)\) of \cref{obj:synth_7}, with the observable marks and zero-based record indexing of \cref{obj:synth_14}. Define \(\widehat F_i(x)\) by clipping this histogram value with the following lower and upper cutoffs, respectively:
\[
\begin{cases}
(c_f,C_f), & i=0,\\
(c_f/4,3C_f/4), & i\in\{1,2\},\\
(0,3C_f/4), & i\in\{3,4,5,6\}.
\end{cases}
\]
\end{definitionv}

\begin{definitionv}{synth_9}[Marked histogram residual]
The residual \(R_{i,K}^{(b)}(x)\) is defined by
\[
R_{i,K}^{(b)}(x)
=
H_{i,K}^{(b)}(x)-\widehat F_i(x),
\]
where the marked histogram \(H_{i,K}^{(b)}\) and the clipped pilot vector \(\widehat F\) are defined in \cref{obj:synth_7,obj:synth_8}, respectively. The parameters are arbitrary nonnegative integers \(n\) and \(K\), data consisting of \(n\) source records and \(n\) target covariates, a coordinate index \(i\) among the seven coordinates, a block index \(b\in\{0,1,2,3\}\), an evaluation point \(x\in\mathbb R\), and pilot-clipping parameters \(c_f,C_f\in\mathbb R\).
\end{definitionv}

\begin{definitionv}{synth_6}[Cubic correction resolution]
The cubic resolution \(K_3(n)\), for \(n\in\mathbb N\), is defined by
\[
K_3(n)=
\begin{cases}
1, & n<n_0,\\[2pt]
2^{\left\lfloor \dfrac{\log m_0}{\log 2}\right\rfloor}, & n\ge n_0,
\end{cases}
\]
where \(n_0\) is the threshold for using the cubic estimator and \(m_0\) is the size of block \(0\) at sample size \(n\). The floor is taken in the nonnegative integers, with value \(0\) when its argument is negative; the logarithm at \(0\) is assigned value \(0\).
\end{definitionv}

\begin{definitionv}{synth_5}[Quadratic correction resolution]
The quadratic resolution \(K_2(n)\), for \(n\in\mathbb N\), is defined by
\[
K_2(n)=
\begin{cases}
1, & n<n_0,\\[2pt]
2^{\left\lfloor \frac{4}{3}\,\frac{\log m_0}{\log 2}\right\rfloor},
& n\ge n_0,
\end{cases}
\]
where \(n_0\) is the threshold and \(m_0\) is the size of block \(0\) at sample size \(n\). The floor is taken in the nonnegative integers, with value \(0\) when its argument is negative.
\end{definitionv}

\begin{algorithmv}{def:cubic-estimator}[Cubic estimator \(\widehat T_A\)]
For \(A\in\{Y,D\}\), the inputs are the sample size \(n\), the threshold \(n_0\), the deterministic sample-channel blocks \(\mathcal B_b^G\) of sizes \(m_b\), the dyadic resolution maps \(K_0(n),K_2(n),K_3(n)\), the normalized marked histograms \(H_{i,K}^{(b)}\), the clipped pilot vector \(\widehat F\), the histogram residuals \(R_{i,K}^{(b)}\), the seven-coordinate index set \(\mathcal I\), the coordinate channels \(g(i)\), the observable marks \(W_{i,r}\), and the observed covariates \(X_r^{g(i)}\).
\begin{enumerate}
\item For \(n\ge n_0\), set the resolutions and cell midpoints to
\[
K_0=K_0(n),\qquad K_2=K_2(n),\qquad K_3=K_3(n),
\qquad
x_{\ell,K}=\frac{\ell-\tfrac12}{K},
\quad \ell=1,\ldots,K.
\]
The resolutions satisfy \(K_0\mid K_3\) and \(K_3\mid K_2\).

\item For \(n\ge n_0\), form the observable linear correction
\[
\mathsf L_A=
\sum_{i\in\mathcal I}
\left\{
\frac{1}{m_1}
\sum_{r\in\mathcal B_1^{g(i)}}
W_{i,r}\,\partial_i\Phi_A
\bigl(\widehat F(X_r^{g(i)})\bigr)
-
\int_0^1
\partial_i\Phi_A(\widehat F(x))\widehat F_i(x)\,dx
\right\}.
\]

\item For \(n\ge n_0\), form the observable quadratic correction
\[
\mathsf Q_A=
\frac{1}{2K_2}
\sum_{i,j\in\mathcal I}\sum_{\ell=1}^{K_2}
\partial_{ij}\Phi_A(\widehat F(x_{\ell,K_2}))
R_{i,K_2}^{(1)}(x_{\ell,K_2})
R_{j,K_2}^{(2)}(x_{\ell,K_2}).
\]

\item For \(n\ge n_0\), form the observable cubic correction
\[
\mathsf C_A=
\frac{1}{6K_3}
\sum_{i,j,k\in\mathcal I}\sum_{\ell=1}^{K_3}
\partial_{ijk}\Phi_A(\widehat F(x_{\ell,K_3}))
R_{i,K_3}^{(1)}(x_{\ell,K_3})
R_{j,K_3}^{(2)}(x_{\ell,K_3})
R_{k,K_3}^{(3)}(x_{\ell,K_3}).
\]
Repeated coordinate indices use distinct evaluation blocks.

\item Output
\[
\widehat T_A=
\begin{cases}
\displaystyle
\int_0^1\Phi_A(\widehat F(x))\,dx+
\mathsf L_A+\mathsf Q_A+\mathsf C_A,
& n\ge n_0,\\
0,& n<n_0.
\end{cases}
\]
Every integral is an exact finite sum, with \(K_0\mid K_3\) and \(K_3\mid K_2\). For \(n<n_0\), the score interval takes the finite-size fallback
\[
I_n=\Theta.
\]
\end{enumerate}
\end{algorithmv}

\begin{lemmav}{lem:observable-marked-reduction}[Observable marked density reduction]
Let \(n\in\mathbb N\) and let \(P\) be a transport law. Suppose:
\begin{itemize}
\item \textbf{(Lower density envelope.)} \(0<c_f<1\).
\item \textbf{(Upper density envelope.)} \(1<C_f\).
\item \textbf{(Hölder radius.)} \(1<L\).
\item \textbf{(Model membership.)} \(P\in\mathcal M_n\) for these constants, as specified in \cref{obj:def:model-class}.
\end{itemize}
Let \(F\), the seven-coordinate marked density vector, be as in \cref{obj:def:marked-density-vector}, with coordinates indexed by \(0,\ldots,6\). Write \(P_{\mathrm T}\) for the target covariate law and \(P_{\mathrm S}\) for the observed source-record law, whose record consists of covariate \(X\), encouragement \(Z\), receipt \(D\), and outcome \(Y\).

For every measurable \(B\subseteq[0,1]\) and every \(i\in\{0,\ldots,6\}\),
\[
\int_B F_i(x)\,dx
=
\begin{cases}
P_{\mathrm T}(B), & i=0,\\
\displaystyle\int_{\{X\in B\}}(1-Z)\,dP_{\mathrm S}, & i=1,\\
\displaystyle\int_{\{X\in B\}}Z\,dP_{\mathrm S}, & i=2,\\
\displaystyle\int_{\{X\in B\}}(1-Z)Y\,dP_{\mathrm S}, & i=3,\\
\displaystyle\int_{\{X\in B\}}ZY\,dP_{\mathrm S}, & i=4,\\
\displaystyle\int_{\{X\in B\}}(1-Z)D\,dP_{\mathrm S}, & i=5,\\
\displaystyle\int_{\{X\in B\}}ZD\,dP_{\mathrm S}, & i=6.
\end{cases}
\]
Moreover, the source arm densities \(q_0\) and \(q_1\) satisfy
\[
q_0(x)=F_1(x)\ge \frac{c_f}{4},
\qquad
q_1(x)=F_2(x)\ge \frac{c_f}{4},
\qquad x\in[0,1].
\]
For each \(A\in\{Y,D\}\), the transported contrast \(T_A\) and the integrand \(\Phi_A\) of \cref{obj:def:smooth-functional} satisfy
\[
T_A=\int_0^1\Phi_A(F(x))\,dx.
\]

For every \(A\in\{Y,D\}\) and every \(v\in\mathbb R^7\) in the deterministic clipping rectangle
\[
c_f\le v_0\le C_f,\qquad
\frac{c_f}{4}\le v_1,v_2\le\frac{3C_f}{4},\qquad
0\le v_i\le\frac{3C_f}{4}\quad(i=3,\ldots,6),
\]
every order \(k\in\{0,\ldots,4\}\) and every choice of coordinate indices \(i_1,\ldots,i_k\in\{0,\ldots,6\}\) obey
\[
\left|
\frac{\partial^k\Phi_A}{\partial v_{i_1}\cdots\partial v_{i_k}}(v)
\right|
\le
48(1+C_f)^2\left(\frac{4}{c_f}\right)^5.
\]
At order zero, the expression on the left is \(\lvert\Phi_A(v)\rvert\).
\end{lemmav}

\begin{lemmav}{lem:marked-cubic-mse}[Uniform cubic mean-square bound]
Let \(c_f\) and \(C_f\) be the lower and upper density envelopes, and let \(L\) be the common Hölder radius. Suppose that:
\begin{itemize}
\item \textbf{(Envelope constants.)} \(0<c_f<1\), \(1<C_f\), and \(1<L\).
\item \textbf{(Sample size.)} The common source and target sample size \(n\) satisfies
\[
n\ge n_0=256.
\]
\item \textbf{(Model.)} \(P\in\mathcal M_n\), with \(\mathcal M_n\) defined in \cref{obj:def:model-class}.
\end{itemize}
For each \(A\in\{Y,D\}\), let \(\widehat T_A\) be the statistic in \cref{obj:def:cubic-estimator}, and let \(T_A\) be the target-averaged source contrast,
\[
T_A
=
\int_{\text{covariate space}}
f_{\mathrm T}(x)\bigl(m_{A1}(x)-m_{A0}(x)\bigr)\,dx.
\]
Then, under the equal-size source–target sampling law specified by \(P\),
\[
\mathbb E_P\!\left[(\widehat T_A-T_A)^2\right]
\le C_{\mathrm{mse}}\,n^{-2/3},
\]
where the finite, explicitly computable constant \(C_{\mathrm{mse}}=C_{\mathrm{mse}}(c_f,C_f,L)\) is
\[
\begin{aligned}
C_{\mathrm{mse}}
={}&3\Bigg\{
\left[2(1+C_f)^2(4/c_f)^5\right]^2 7^8
\Big[
2^7\,8!\,
\big\{(1+C_f+C_f^2)^4 5^{4/5}
 +(1+C_f+C_f^2)5^{7/5}\big\}\\
&\hspace{48mm}
+2^8[3(1+C_f)L]^8 5^{4/5}
\Big]\\
&\quad
+7^4[48(1+C_f)^2(4/c_f)^5]^2
 [3(1+C_f)L]^4 2^{1/2}5^{2/3}\\
&\quad
+2\Big[
\big\{24(1+C_f)^2(4/c_f)^5
       7^2[3(1+C_f)L]^2\big\}^2
 7^2\\
&\hspace{24mm}\cdot
\big\{2(1+C_f+C_f^2)5^{1/5}
 +2^{5/4}[3(1+C_f)L]^2 5^{1/5}\big\}
 2^{1/2}5^{1/2}\\
&\hspace{16mm}
+\big\{8(1+C_f)^2(4/c_f)^5
       7^3[3(1+C_f)L]^3\big\}^2
 2^{3/4}5^{3/4}
\Big]
\Bigg\}\\
&+3\Bigg\{
10\cdot7^2[48(1+C_f)^2(4/c_f)^5]^2\\
&\quad
+60\cdot7^4(2^2-1)^2
 [48(1+C_f)^2(4/c_f)^5]^2
 [1+2(1+C_f)]^4[2+C_f+C_f^2]^2\\
&\quad
+310\cdot7^6(2^3-1)^2
 [48(1+C_f)^2(4/c_f)^5]^2
 [1+2(1+C_f)]^6[2+C_f+C_f^2]^3
\Bigg\}.
\end{aligned}
\]
\end{lemmav}

\begin{theoremv}{thm:honest-upper-full-elbow}[Honest upper length envelope]
Fix the noncoverage probability \(\alpha\), the density envelopes \(c_f,C_f\), and the Hölder radius \(L\), satisfying:
\begin{itemize}
\item \textbf{(Noncoverage.)} \(0<\alpha<1\).
\item \textbf{(Lower density envelope.)} \(0<c_f<1\).
\item \textbf{(Upper density envelope.)} \(1<C_f\).
\item \textbf{(Hölder radius.)} \(1<L\).
\end{itemize}
The score interval \(I_n\) of \cref{obj:def:score-interval} belongs to the class \(\mathfrak H_n\) of honest intervals in \cref{obj:def:honest-intervals} for every positive integer \(n\).

Set the sample-size threshold to
\[
n_0=256.
\]
Write \(\mu(P)\) for the transported first stage, namely the target-averaged source receipt contrast:
\[
\mu(P)=T_D(P).
\]
There exists a finite constant \(C_0=C_0(\alpha,c_f,C_f,L)>0\) such that, for every integer \(n\ge n_0\) and every \(P\in\mathcal M_n\), the model class of \cref{obj:def:model-class},
\[
\mathbb E_P\!\left[\lambda_{\Theta}(I_n)\right]
\le C_0\min\left\{1,\frac{n^{-1/3}}{\mu(P)}\right\},
\]
where \(\lambda_{\Theta}\) denotes Lebesgue length restricted to the effect domain \(\Theta\), and the expectation is under the joint law of the source and target samples, each of size \(n\).

With the same constant \(C_0\), for every integer \(n\ge n_0\) and every real \(a\) satisfying \(0<a\le 1/4\), the minimax expected-length frontier \(L_n(a)\) of \cref{obj:def:length-frontier} satisfies
\[
L_n(a)\le C_0\min\left\{1,\frac{n^{-1/3}}{a}\right\}.
\]
\end{theoremv}

\begin{theoremv}{thm:honest-lower-full-elbow}[Honest length lower bound]
Fix constants satisfying:
\begin{itemize}
\item \textbf{(Coverage.)} The noncoverage probability satisfies \(0<\alpha<1\).
\item \textbf{(Density envelopes.)} The lower and upper density envelopes satisfy \(0<c_f<1<C_f\).
\item \textbf{(Smoothness.)} The common Hölder radius satisfies \(L>1\).
\end{itemize}
Let \(\mathcal M_n\) be the model class in \cref{obj:def:model-class}, and let \(\mathcal M_n(a)\) be its strength slice in \cref{obj:def:strength-slice}. Write \(\theta(P)\) for the target-population average treatment effect among compliers under \(P\), and let the effect domain and sample-size threshold be
\[
\Theta=[-1,1],
\qquad n_0=256.
\]
There exists a constant \(c_0=c_0(\alpha,c_f,C_f,L)>0\) such that, for every integer \(n\ge n_0\) and every \(0<a\le 1/4\), the following holds. Every measurable confidence-set procedure \(C_n\), based on \(n\) source records and \(n\) target covariate records, with \(C_n(\omega)\subseteq\Theta\) for every sample realization \(\omega\), and satisfying
\[
P\bigl(\theta(P)\in C_n\bigr)\ge 1-\alpha
\qquad\text{for every }P\in\mathcal M_n,
\]
obeys
\[
\sup_{P\in\mathcal M_n(a)}
\mathbb E_P\!\left[\lambda_{\Theta}(C_n)\right]
\ge c_0\min\left\{1,\frac{n^{-1/3}}{a}\right\},
\qquad
\lambda_{\Theta}(C_n)
=\operatorname{Leb}(C_n\cap\Theta).
\]
Here probabilities and expectations are taken under the observed two-sample law. Moreover, the minimax expected-length frontier \(L_n(a)\) of \cref{obj:def:length-frontier} satisfies
\[
L_n(a)\ge c_0\min\left\{1,\frac{n^{-1/3}}{a}\right\}.
\]
\end{theoremv}

\begin{definitionv}{synth_18}[Resolution-\(K\) interval cells]
For a positive integer \(K\) and an index \(\ell\in\{0,\ldots,K-1\}\), define
\[
\operatorname{cell}(K,\ell)=
\begin{cases}
[\ell/K,(\ell+1)/K),&0\le\ell<K-1,\\
[(K-1)/K,1],&\ell=K-1.
\end{cases}
\]
Thus each cell includes its left endpoint, and the final cell also includes the endpoint \(1\).
\end{definitionv}

\begin{definitionv}{synth_12}[Tiled and coupled perturbations]
The tiled perturbation \(u_{\lambda}\), for \(c_{\star}\in\mathbb R\), \(n\in\mathbb N\), and a cellwise sign vector \(\lambda\in\{-1,1\}^{K_{\mathrm L}}\), is defined for every \(x\in\mathbb R\) by
\[
K_{\mathrm L}=\left\lceil n^{4/3}\right\rceil,
\qquad
h_n=c_{\star}K_{\mathrm L}^{-1/8},
\qquad
u_{\lambda}(x)=
\sum_{\ell=0}^{K_{\mathrm L}-1}
\begin{cases}
h_n\lambda_\ell\sin\!\bigl(2\pi(K_{\mathrm L}x-\ell)\bigr),
& x\in\operatorname{cell}(K_{\mathrm L},\ell),\\
0,&\text{otherwise}.
\end{cases}
\]
Each sign \(\lambda_\ell\) equals \(1\) when the supplied Boolean entry for cell \(\ell\) is true and \(-1\) when it is false. For \(0\leq\ell<K_{\mathrm L}\), the cells are
\[
\operatorname{cell}(K_{\mathrm L},\ell)=
\begin{cases}
[\ell/K_{\mathrm L},1],
& \ell+1=K_{\mathrm L},\\
[\ell/K_{\mathrm L},(\ell+1)/K_{\mathrm L}),
& \ell+1\ne K_{\mathrm L}.
\end{cases}
\]
At \(n=0\), the sum is empty and equals zero, with the convention \(0^{-1/8}=0\).
The coupled perturbation \(v_{\lambda,\tau}\), for any \(\tau\in\mathbb R\), is defined by
\[
v_{\lambda,\tau}(x)=\tau u_{\lambda}(x),
\qquad x\in\mathbb R.
\]
\end{definitionv}

\begin{definitionv}{synth_15}[Baseline receipt–outcome cell laws]
For a baseline compliance share \(b\in(0,1/4]\) and \(z\in\{0,1\}\), set
\[
r_z^{(b)}=\frac{1+(2z-1)b}{2}.
\]
Define the baseline conditional receipt–outcome probability mass function on \(\{0,1\}^2\) by
\[
R_z^{(b)}(d,y)=
\begin{cases}
r_z^{(b)}/2,&d=1,\\
(1-r_z^{(b)})/2,&d=0,
\end{cases}
\qquad z,d,y\in\{0,1\}.
\]
Equivalently, for each \(y\in\{0,1\}\),
\[
R_0^{(b)}(0,y)=R_1^{(b)}(1,y)=\frac{1+b}{4},
\qquad
R_0^{(b)}(1,y)=R_1^{(b)}(0,y)=\frac{1-b}{4}.
\]
Under this baseline law, receipt has probability \(r_z^{(b)}\) conditional on instrument value \(z\), and the binary outcome is independent of receipt with probability \(1/2\) for each outcome value. In the lower-mixture construction, the baseline share is \(b=\max\{a,n^{-1/3}\}\) for \(n\ge n_0=256\) and \(0<a\le1/4\).
\end{definitionv}

\begin{algorithmv}{def:legal-iv-mixture}[Legal IV mixture construction]
The legal IV mixture is the triple \((Q_{\lambda,\tau},M_{\tau},P_{\star})\), consisting of a component law, its uniform sign-mixture experiment, and a center law, defined for the following parameters:
\begin{itemize}
\item \textbf{(Sample size.)} \(n\in\mathbb N\) satisfies
\[
n\ge n_0=256,
\qquad
K_{\mathrm L}=\left\lceil n^{4/3}\right\rceil.
\]
Here \(K_{\mathrm L}\) is the number of perturbation cells.
\item \textbf{(Baseline strength.)} The strength parameter \(a\) satisfies \(0<a\le1/4\), and the baseline compliance share is
\[
b=\bar a=\max\{a,n^{-1/3}\}.
\]
\item \textbf{(Amplitude and admissibility.)} The amplitude \(c_{\star}\) satisfies \(c_{\star}<1\). With perturbation height
\[
h=h_n=c_{\star}K_{\mathrm L}^{-1/8},
\]
require \(n>0\) and \(\operatorname{Adm}_{n,a}(c_{\star})\), where
\[
\operatorname{Adm}_{n,a}(c_{\star})
\Longleftrightarrow
\left(
0<c_{\star}
\ \text{and}\
h\le\frac{3}{16}
\ \text{and}\
h^2\le b\left(\frac14-h^2\right)
\right).
\]
\item \textbf{(Sign and continuum indices.)} The cellwise sign vector and continuum index satisfy
\[
\lambda\in\{-1,1\}^{K_{\mathrm L}},
\qquad
\tau\in[-1,1].
\]
For the tiled covariate perturbation \(u_{\lambda}\) indexed by these signs, set
\[
v_{\lambda,\tau}=\tau u_{\lambda}.
\]
\end{itemize}

Define the source and target covariate densities and encouragement probabilities by
\[
f_{\mathrm S}=1,
\qquad
f_{\mathrm T}=1+u_{\lambda},
\qquad
e_1^{(\lambda)}=\frac12+u_{\lambda},
\qquad
e_0^{(\lambda)}=\frac12-u_{\lambda}.
\]
For \(z,d,y\in\{0,1\}\), the source joint cells conditional on covariates are
\[
Q_{\lambda,\tau}(Z=z,D=d,Y=y\mid X=x)
=
e_z^{(\lambda)}(x)R_z^{(b)}(d,y)
+\frac{v_{\lambda,\tau}(x)(2y-1)}4,
\]
where \(R_z^{(b)}(d,y)\) is the baseline receipt–outcome cell probability at compliance share \(b\). The cells conditional on encouragement and covariates are
\[
Q_{\lambda,\tau}(D=d,Y=y\mid Z=z,X=x)
=
\frac{
e_z^{(\lambda)}(x)R_z^{(b)}(d,y)
+v_{\lambda,\tau}(x)(2y-1)/4
}{
e_z^{(\lambda)}(x)
}.
\]

Let \(C\) be the complier indicator, \(D(z)\) potential receipt under encouragement \(z\), and \(Y(d)\) the potential outcome under receipt state \(d\). Define the complier outcome margins by
\[
M_{1y}(x)
=
P(C=1,Y(1)=y\mid X=x)
=
\frac b2
-
\frac{(2y-1)u_{\lambda}(x)v_{\lambda,\tau}(x)}
{2\{1/4-u_{\lambda}(x)^2\}},
\]
\[
M_{0y}(x)
=
P(C=1,Y(0)=y\mid X=x)
=
\frac b2
+
\frac{(2y-1)u_{\lambda}(x)v_{\lambda,\tau}(x)}
{2\{1/4-u_{\lambda}(x)^2\}},
\qquad y\in\{0,1\}.
\]
Assign always-takers and never-takers the conditional masses
\[
P(D(0)=D(1)=1,Y(1)=y\mid X=x)
=
Q_{\lambda,\tau}(D=1,Y=y\mid Z=0,X=x),
\]
\[
P(D(0)=D(1)=0,Y(0)=y\mid X=x)
=
Q_{\lambda,\tau}(D=0,Y=y\mid Z=1,X=x).
\]
For compliers, set \(D(0)=0\), \(D(1)=1\), and
\[
P(Y(d)=y\mid X=x,C=1)
=
\frac{M_{dy}(x)}b,
\qquad d,y\in\{0,1\},
\]
coupling \(Y(0)\) and \(Y(1)\) independently under these binary marginals. Set
\[
Y(0)=0\quad\text{for always-takers},
\qquad
Y(1)=0\quad\text{for never-takers}.
\]
Draw encouragement according to
\[
Z\ \perp\ (D(0),D(1),Y(0),Y(1))\mid X,
\qquad
P(Z=1\mid X=x)=e_1^{(\lambda)}(x),
\]
and observe
\[
D=D(Z),
\qquad
Y=Y(D).
\]
Use the same conditional primitive full-data law given \(X\) in source and target. Represent the resulting combined law \(Q_{\lambda,\tau}\) by
\[
Q_{\lambda,\tau}(S=\mathrm{source})
=
Q_{\lambda,\tau}(S=\mathrm{target})
=
\frac12.
\]
This allocation is a convention for representing the constructed witness; admissibility specifies the parameter domain of the construction.

Define the center law by
\[
P_{\star}
=
\left.Q_{\lambda,\tau}\right|_{u_{\lambda}=v_{\lambda,\tau}=0},
\]
retaining the baseline compliance share \(b\) and the representation convention
\[
P_{\star}(S=\mathrm{source})
=
P_{\star}(S=\mathrm{target})
=
\frac12.
\]
For each component and for the center, take independent source and target product samples, each of size \(n\). Writing \(P_{\mathrm S}\) and \(P_{\mathrm T}\) for the one-record observed source and target laws induced by the law in question, the center experiment and uniform sign mixture are
\[
P_{\star}^{(n,n)}
=
\left.
\left(P_{\mathrm S}^{\otimes n}\otimes P_{\mathrm T}^{\otimes n}\right)
\right|_{P_{\star}},
\]
\[
M_{\tau}
=
2^{-K_{\mathrm L}}
\sum_{\lambda\in\{-1,1\}^{K_{\mathrm L}}}
\left.
\left(P_{\mathrm S}^{\otimes n}\otimes P_{\mathrm T}^{\otimes n}\right)
\right|_{Q_{\lambda,\tau}}.
\]
\end{algorithmv}

\begin{definitionv}{synth_17}[Distances between observed experiments]
For probability measures \(P\) and \(Q\) on the same measurable observation space \((\Omega,\mathcal F)\), define the chi-square divergence by
\[
\chi^2(P,Q)=
\begin{cases}
\displaystyle\int_\Omega
\left(\frac{dP}{dQ}-1\right)^2\,dQ,&P\ll Q,\\
+\infty,&P\not\ll Q.
\end{cases}
\]
The integral is understood in the extended nonnegative reals. Define total-variation distance with the probability-distance normalization by
\[
\operatorname{TV}(P,Q)
=\sup_{A\in\mathcal F}|P(A)-Q(A)|
=\frac12\int_\Omega
\left|\frac{dP}{d\nu}-\frac{dQ}{d\nu}\right|\,d\nu,
\qquad \nu=P+Q.
\]
Thus \(\operatorname{TV}(P,Q)\in[0,1]\). For the lower mixture, take \(P=M_\tau\) and \(Q=P_\star^{(n,n)}\) on the observation space of the independent source and target samples, each of size \(n\).
\end{definitionv}

\begin{lemmav}{lem:legal-iv-mixture-full-elbow}[Legal mixture and effect separation]
Fix the prescribed noncoverage probability \(\alpha\), density envelopes \(c_f,C_f\), and Hölder radius \(L\), satisfying:
\begin{itemize}
\item \textbf{(Noncoverage.)} \(0<\alpha<1\).
\item \textbf{(Lower envelope.)} \(0<c_f<1\).
\item \textbf{(Upper envelope.)} \(1<C_f\).
\item \textbf{(Radius.)} \(1<L\).
\end{itemize}
There exists an amplitude \(c_{\star}\in(0,1)\), chosen uniformly in the sample size and strength floor, with the following properties. For every integer \(n\ge n_0=256\) and every \(0<a\le1/4\), set
\[
K_{\mathrm L}=\lceil n^{4/3}\rceil,
\qquad
\bar a=\max\{a,n^{-1/3}\},
\qquad
h_n=c_{\star}K_{\mathrm L}^{-1/8}.
\]
The amplitude satisfies the construction's admissibility conditions
\[
c_{\star}>0,
\qquad
h_n\le\frac{3}{16},
\qquad
h_n^2\le\bar a\left(\frac14-h_n^2\right).
\]

For every \(\tau\in[-1,1]\) and every cellwise sign vector
\(\lambda\in\{-1,1\}^{K_{\mathrm L}}\), the constructed component
\(Q_{\lambda,\tau}\) belongs to the strength slice
\(\mathcal M_n(a)\) of \cref{obj:def:strength-slice}.
For each \(s\in\{0,1\}\) and every \(x\in[0,1]\), its pointwise conditional-mass averages satisfy
\[
\mathbb E_{Q_{\lambda,\tau}}[C\mid X=x,S=s]=\bar a,
\qquad
\mathbb E_{Q_{\lambda,\tau}}
 [(Y(1)-Y(0))C\mid X=x,S=s]
 =
 -\frac{u_{\lambda}(x)v_{\lambda,\tau}(x)}
 {1/4-u_{\lambda}(x)^2},
\]
where \(C=\mathbf 1\{D(1)>D(0)\}\), \(u_{\lambda}\) is the construction's tiled covariate perturbation, and \(v_{\lambda,\tau}=\tau u_{\lambda}\). These pointwise averages also define versions of the corresponding conditional means under \(Q_{\lambda,\tau}\). Its transported first stage \(\mu\) and target complier effect \(\theta\) satisfy
\[
\mu(Q_{\lambda,\tau})=\bar a,
\qquad
\theta(Q_{\lambda,\tau})=-\frac{\tau J_n}{\bar a},
\]
where the common separation functional \(J_n\) is characterized by
\[
J_n=-\bar a\,\theta(Q_{\lambda,1}).
\]

The unperturbed center \(P_{\star}\) belongs to
\(\mathcal M_n(a)\), and
\[
\mu(P_{\star})=\bar a,
\qquad
\theta(P_{\star})=0,
\qquad
2^{3/4}c_{\star}^2n^{-1/3}
\le J_n\le4c_{\star}^2n^{-1/3}.
\]
For the uniform sign mixture and the center's equal-size observed product experiment,
\[
M_{\tau}
=
2^{-K_{\mathrm L}}
\sum_{\lambda\in\{-1,1\}^{K_{\mathrm L}}}
\left(Q_{\lambda,\tau,\mathrm S}^{\otimes n}
\otimes Q_{\lambda,\tau,\mathrm T}^{\otimes n}\right),
\qquad
P_{\star}^{(n,n)}
=
P_{\star,\mathrm S}^{\otimes n}
\otimes P_{\star,\mathrm T}^{\otimes n},
\]
the subscripts \(\mathrm S,\mathrm T\) denote the observed source-record and target-covariate laws. Uniformly over \(\tau\in[-1,1]\),
\[
1+\chi^2(M_{\tau},P_{\star}^{(n,n)})
\le
\exp\{C_{\mathrm{mix}}c_{\star}^4\},
\qquad
C_{\mathrm{mix}}=\frac{19321}{1800},
\qquad
\operatorname{TV}(M_{\tau},P_{\star}^{(n,n)})
\le\frac{1-\alpha}{2}.
\]
Moreover, \(M_{\tau}\ll P_{\star}^{(n,n)}\), and its likelihood ratio minus one is square-integrable under \(P_{\star}^{(n,n)}\).
\end{lemmav}

\begin{definitionv}{synth_16}[Tiled perturbation bump]
Define \(\operatorname{bump}:\mathbb R\to[-1,1]\) by
\[
\operatorname{bump}(t)=
\begin{cases}
\sin(2\pi t),&t\in[0,1],\\
0,&t\notin[0,1].
\end{cases}
\]
Its restriction to \([0,1]\) is the function \(B(t)=\sin(2\pi t)\) used in the tiled perturbation. Its support is \([0,1]\), and its normalization is
\[
\|\operatorname{bump}\|_\infty=1,\qquad
\int_0^1\operatorname{bump}(t)\,dt=0,\qquad
\int_0^1\operatorname{bump}(t)^2\,dt=\frac12.
\]
It satisfies \(\operatorname{bump}(1-t)=-\operatorname{bump}(t)\) for \(t\in[0,1]\) and vanishes at both endpoints. The zero extension is continuous and globally Lipschitz with constant \(2\pi\); its restriction to \([0,1]\) is smooth. In particular, for \(\beta=1/8\),
\[
|\operatorname{bump}(s)-\operatorname{bump}(t)|
\le 2\pi |s-t|^\beta,\qquad s,t\in\mathbb R.
\]
\end{definitionv}

\begin{lemmav}{lem:random-set-length-identity}[Expected random-set length]
Let \(C_n\subseteq\Theta\) be a measurable random set under \(P\), where \(\Theta\) is the effect domain and \(\lambda_{\Theta}\) denotes Lebesgue length restricted to that domain. Then
\[
\mathbb E_P\!\left[\lambda_{\Theta}(C_n)\right]
=
\int_{\Theta}P(t\in C_n)\,dt.
\]
\end{lemmav}

\begin{definitionv}{def:fixed-geometry-subclass}[Fixed-geometry subclass $\mathcal M_n^{\mathrm{fix}}(\mathcal G)$]
The fixed-geometry subclass is
\[
\mathcal M_n^{\mathrm{fix}}(\mathcal G)
=
\bigl\{P\in\mathcal M_n:(f_{\mathrm S},f_{\mathrm T},e)=\mathcal G\bigr\},
\]
with equality at sample-size index \(n\). Here \(\mathcal G\) is the supplied triple of source covariate density \(f_{\mathrm S}\), target covariate density \(f_{\mathrm T}\), and source encouragement propensity \(e\). Membership in \(\mathcal M_n\) requires all of the following restrictions at that index:
\begin{itemize}
\item \textbf{(Sampling.)} \cref{obj:ass:source-iid,obj:ass:target-iid,obj:ass:sample-independence}.
\item \textbf{(Covariate densities.)} \cref{obj:ass:source-density-bounds,obj:ass:target-density-bounds,obj:ass:source-density-holder,obj:ass:target-density-holder}.
\item \textbf{(Propensity and arm means.)} \cref{obj:ass:propensity-overlap,obj:ass:propensity-holder,obj:ass:arm-means-holder}.
\item \textbf{(Causal restrictions.)} \cref{obj:ass:receipt-consistency,obj:ass:outcome-consistency,obj:ass:instrument-randomization,obj:ass:exclusion,obj:ass:monotonicity}.
\item \textbf{(Transport and strength.)} \cref{obj:ass:transport-complier-outcome,obj:ass:transport-complier-share,obj:ass:positive-strength}.
\end{itemize}
\end{definitionv}

\begin{definitionv}{def:known-geometry-score}[Supplied-geometry score $\widetilde T_A^{\mathcal G}$]
For \(A\in\{Y,D\}\) and every positive integer \(n\), the supplied-geometry score is
\[
\widetilde T_A^{\mathcal G}
=
\frac1n\sum_{i=1}^n
\frac{f_{\mathrm T}^0(X_i)}{f_{\mathrm S}^0(X_i)}
\left\{
\frac{Z_iA_i}{e^0(X_i)}
-
\frac{(1-Z_i)A_i}{1-e^0(X_i)}
\right\}.
\]
Here \((X_i,Z_i,D_i,Y_i)\) are the observed source records, comprising covariates, encouragement, receipt, and outcome, and the supplied geometry is
\[
\mathcal G=(f_{\mathrm S}^0,f_{\mathrm T}^0,e^0),
\]
where \(f_{\mathrm S}^0\) and \(f_{\mathrm T}^0\) are the supplied source and target covariate densities and \(e^0\) is the supplied source encouragement propensity. The statistic is a function of the observed source sample and \(\mathcal G\) alone.
\end{definitionv}

\begin{definitionv}{synth_11}[Known geometry calibration radius]
The supplied-geometry calibration radius is defined by
\[
\rho_n^{\mathcal G}
=
\frac{4C_f}{c_f}
\sqrt{\frac{2\log(4/\alpha)}{n}},
\]
for real parameters \(\alpha\), the prescribed noncoverage probability, \(C_f\), the upper density envelope, and \(c_f\), the lower density envelope, and a natural-number sample size \(n\), satisfying:
\begin{itemize}
\item \textbf{(Noncoverage probability.)} \(0<\alpha<1\).
\item \textbf{(Upper envelope.)} \(1<C_f\).
\item \textbf{(Lower envelope.)} \(0<c_f<1\).
\item \textbf{(Sample size.)} \(0<n\).
\end{itemize}
\end{definitionv}

\begin{definitionv}{def:known-geometry-interval}[Supplied-geometry interval $I_n^{\mathcal G}$]
The supplied-geometry interval is
\[
I_n^{\mathcal G}
=
\begin{cases}
\mathcal A_n^{\mathcal G},&\mathcal A_n^{\mathcal G}\ne\varnothing,\\
\{0\},&\mathcal A_n^{\mathcal G}=\varnothing,
\end{cases}
\]
where the affine-score acceptance set is
\[
\mathcal A_n^{\mathcal G}
=
\left\{
v\in\Theta:
\left|\widetilde T_Y^{\mathcal G}
-v\widetilde T_D^{\mathcal G}\right|
\le 2\rho_n^{\mathcal G}
\right\}.
\]
Here \(\Theta\) is the bounded domain of the complier effect, \(\rho_n^{\mathcal G}\) is the calibration radius, and \(\widetilde T_Y^{\mathcal G}\) and \(\widetilde T_D^{\mathcal G}\) are the supplied-geometry outcome and receipt scores.
\end{definitionv}

\begin{definitionv}{def:finite-cell-subclass}[Finite-cell subclass $\mathcal M_n^{\mathrm{cell}}(\Pi)$]
The finite-cell subclass \(\mathcal M_n^{\mathrm{cell}}(\Pi)\) consists of the laws \(P\in\mathcal M_n\) whose source covariate density \(f_{\mathrm S}\), target covariate density \(f_{\mathrm T}\), and source encouragement propensity \(e\) are constant on every cell of the fixed finite interval partition \(\Pi\), at sample-size index \(n\). Membership in \(\mathcal M_n\) requires all of the following restrictions at that index:
\begin{itemize}
\item \textbf{(Sampling.)} \cref{obj:ass:source-iid,obj:ass:target-iid,obj:ass:sample-independence}.
\item \textbf{(Covariate densities.)} \cref{obj:ass:source-density-bounds,obj:ass:target-density-bounds,obj:ass:source-density-holder,obj:ass:target-density-holder}.
\item \textbf{(Propensity and arm means.)} \cref{obj:ass:propensity-overlap,obj:ass:propensity-holder,obj:ass:arm-means-holder}.
\item \textbf{(Causal restrictions.)} \cref{obj:ass:receipt-consistency,obj:ass:outcome-consistency,obj:ass:instrument-randomization,obj:ass:exclusion,obj:ass:monotonicity}.
\item \textbf{(Transport and strength.)} \cref{obj:ass:transport-complier-outcome,obj:ass:transport-complier-share,obj:ass:positive-strength}.
\end{itemize}
\end{definitionv}

\begin{definitionv}{synth_13}[Centered covariate witness direction]
The witness direction \(g_{\dagger}\) is defined, for every \(x\in\mathbb R\), by
\[
g_{\dagger}(x)
=
\left|x-\frac12\right|^{1/8}
-
\int_0^1 \left|u-\frac12\right|^{1/8}\,du.
\]
\end{definitionv}

\begin{definitionv}{def:continuous-geometry-witness}[Continuous-geometry laws $P_{\sigma}^{\dagger}$]
For each \(\sigma\in\{-1,1\}\), the full-data law \(P_{\sigma}^{\dagger}\) is specified by the following covariate densities, principal-stratum distribution, potential variables, and sampling scheme. Define the perturbation amplitude by
\[
\varepsilon_{\dagger}
=
\frac14\min\{1-c_f,C_f-1,1/4,L-1\},
\]
where \(c_f\) and \(C_f\) are the fixed lower and upper density envelopes and \(L\) is the common Hölder radius. With \(g_{\dagger}\) denoting the covariate perturbation, set
\[
f_{\mathrm S}(x)=1+\sigma\varepsilon_{\dagger}g_{\dagger}(x),
\qquad
f_{\mathrm T}(x)=1-\sigma\varepsilon_{\dagger}g_{\dagger}(x),
\qquad
e(x)=\frac12+\sigma\varepsilon_{\dagger}g_{\dagger}(x).
\]
The law is completed as follows:
\begin{itemize}
\item \textbf{(Principal strata.)} Conditional on \(X\) in each population, the complier, always-taker, and never-taker probabilities are, respectively,
\[
\frac18,\qquad \frac7{16},\qquad \frac7{16}.
\]
\item \textbf{(Potential variables.)} The potential receipts \(D(0)\) and \(D(1)\) are the receipts corresponding to the assigned principal stratum, and the potential outcomes satisfy
\[
Y(0)=Y(1)=\frac12.
\]
\item \textbf{(Source observations.)} Draw encouragement \(Z\) conditionally independently of the potential variables given \(X\), with propensity \(e(X)\), and observe
\[
D=D(Z),\qquad Y=Y(D).
\]
\item \textbf{(Samples.)} Generate independent source and target samples with covariate densities \(f_{\mathrm S}\) and \(f_{\mathrm T}\), respectively.
\end{itemize}
For the combined full-data representation of these two conditional population laws, the population indicator \(S\) has probabilities
\[
P_{\sigma}^{\dagger}(S=\mathrm{source})
=
P_{\sigma}^{\dagger}(S=\mathrm{target})
=
\frac12.
\]
These probabilities specify the representation convention for \(P_{\sigma}^{\dagger}\); the ambient model class retains its stated assumptions.
\end{definitionv}
